\medskip

Soit $g$ la fonction définie sur l'intervalle $]0~;~+\infty[$ par 

\[g(x) = 1+ x^2[1 - 2 \ln (x)].\]

La fonction $g$ est dérivable sur l'intervalle $]0~;~+\infty[$ et on note $g'$ sa fonction dérivée.

On appelle $\mathcal{C}$ la courbe représentative de la fonction $g$ dans un repère orthonormé du plan.

\medskip

\textbf{PARTIE A}

\medskip

\begin{enumerate}
\item Justifier que $g(\text{e})$ est strictement négatif.
\item Justifier que $\displaystyle\lim_{x \to + \infty} g(x) = - \infty$.
\item 
	\begin{enumerate}
		\item Montrer que, pour tout $x$ appartenant à l'intervalle $]0~;~+\infty[$,\: $g'(x) = -4x \ln (x)$.
		\item Étudier le sens de variation de la fonction $g$ sur l'intervalle $]0~;~+\infty[$ 
		\item Montrer que l'équation $g(x) = 0$ admet une unique solution, notée $\alpha$, sur l'intervalle $[1~;~+\infty[$.
		\item Donner un encadrement de $\alpha$ d'amplitude $10^{-2}$
	\end{enumerate}
\item Déduire de ce qui précède le signe de la fonction $g$ sur l'intervalle $[1~;~+\infty[$.
\end{enumerate}

\medskip

\textbf{PARTIE B}

\medskip

\begin{enumerate}
\item On admet que, pour tout $x$ appartenant à l'intervalle $[1~;~\alpha]$, $g''(x) = - 4[\ln (x) + 1]$.

Justifier que la fonction $g$ est concave sur l'intervalle $[1~;~\alpha]$.
\end{enumerate}

\begin{minipage}{0.48\linewidth}
\begin{enumerate}[resume]
\item Sur la figure ci-contre, A et B sont les points de la
courbe $\mathcal{C}$ d'abscisses respectives 1 et $\alpha$.
	\begin{enumerate}
		\item Déterminer l'équation réduite de la droite (AB).
		\item En déduire que pour tout réel $x$ appartenant à
l'intervalle $[1~;~\alpha]$,\: $g(x) \geqslant \dfrac{- 2}{\alpha - 1} x + \dfrac{2\alpha}{\alpha - 1}$.
	\end{enumerate}
\end{enumerate}
\end{minipage} \hfill
\begin{minipage}{0.48\linewidth}
\begin{center}
\psset{unit=1.3cm,arrowsize=2pt 3}
\begin{pspicture}(-0.2,-0.2)(2.5,3)
\psgrid[gridlabels=0pt,subgriddiv=1,gridwidth=0.15pt]
\psaxes[linewidth=1.25pt,labelFontSize=\scriptstyle]{->}(0,0)(0,0)(2.5,3)
\psplot[plotpoints=1000,linewidth=1.25pt,linecolor=red]{1}{1.9}{1 x ln 2 mul sub x dup mul mul 1 add}
\uput[ul](1,2){\red A}\uput[ur](1.9,0){\red B}\uput[d](1.9,0.05){$\alpha$}
\uput[ur](1.5,1.5){\red $\mathcal{C}$}
\end{pspicture}
\end{center}
\end{minipage}

\bigskip
